\section{Introduction}
\label{introduction}

\begin{note}
We should focus on the fact that a similar thing has been done before
(\cite{fexpr-shutt}\cite{fexpr-pe}), and that the novel thing we're trying
here is \textbf{callsite-annotated macro expansion}.
Parts of this section should be moved to \cref{discussion}.
\end{note}

We present a programming language that is as minimal as possible
while being expressive enough for general-purpose use. \paraphrase{We are aiming towards} the
following properties:
\begin{enumerate}
  \item Minimality \label{c-minimality}
  \begin{enumerate}
    \item Homoiconicity, provided by LISP-like syntax
    \item Immutability
    \item Few (and simple) special forms
  \end{enumerate}
  \item Expressiveness
  \begin{enumerate}
    \item Mutual recursion
    \item Metaprogramming (allow implementing \paraphrase{syntactic conveniences} as
      libraries written in the language rather than compiler/interpreter
      features)
  \end{enumerate}
  \item Performance \label{c-performance}
\end{enumerate}

Some of these properties are at odds at each other: in particular, it is
difficult to provide mutual recursion and immutability while at the same time
keeping special forms few and simple. For example, Scheme, LISP and other
similar languages forgo the ``immutability'' constraint, which makes it easy to
implement cyclic data structures like mutually-recursive function definitions
\cite[Chapter 4.1.5]{sicp}. The toplevel expressions in such
languages are usually \emph{statements} like \ildmono{define} that \emph{mutate} a global
\emph{environment}.

\begin{schemecode}
(define (even? x)
    (if (= x 0)
        #t
        (odd? (- 1 x))))

(define (odd? x)
    (if (= x 0)
        #t
        (even? (- 1 x))))

(display (even? 42))
\end{schemecode}
\lstcaption{Function definition in Scheme mutates the global environment}

Other LISP-like languages, such as LFE, guarantee immutability of all
data, but handle a lot of the complexity in the interpreter itself: the
language features are written in the host language that implements the
interpreter, and not in the language itself\citneeded. For example, functions defined in
the global namespace are distinct from locally defined lambda objects, and the
interpreter takes special care to support recursion and mutual
recursion without allowing programs to mutate data. In fact, in LFE it is not
even possible to create a cyclic data structure altogether! The price that is
paid to achieve this is that the global namespace of defined functions is not
manipulable by the program (which violates homoiconicity to some
extent) and that \ildmono{define} and similar constructs are special forms.

It is therefore interesting to see if we can design a language that meets all
these goals at the same time. We define a language (which we will call ILD)
which adheres to the \clink{c-minimality}{minimality} constraints, and then
demonstrate expressiveness by implementing increasingly-complex constructs in
ILD, including a way to define mutually-recursive functions (\cref{letrec}).

The techniques applied in this paper are similar to research regarding
F-expressions \cite{fexpr-shutt}\cite{fexpr-pe}. However, instead of declaring
whether an object is a function or macro upon construction, ILD differentiates
between function and macro calls at the callsite (\cref{macroexpand-mechanism}).

ILD tries to \paraphrase{fill the spot} for a LISP-like language that supports
first-class continuations and macros with a minimal core.

\begin{cfig}{Table}{Feature comparison of similar languages}
\begin{tabular}{clccccll}
\toprule
\multicolumn{4}{c}{} & \multicolumn{4}{c}{Macros} \\
\cmidrule(lr){5-8}
Language & Immutable? & Continuations? & Special forms & Kind & First-class? & Expansion & Trigger \\
\midrule
LISP & no & yes & many & unhygienic & no & separate pass & bind time \\
Scheme & no & yes & few & hygienic & no & separate pass & bind time \\
LFE & yes & no & many & limited & no & separate pass & bind time \\
Kraken \cite{fexpr-pe} & yes & no & one & fexprs & yes & runtime & bind time \\
ILD & yes & yes & three & unhygienic & yes & runtime & callsite \\
\bottomrule
\end{tabular}
\end{cfig}

To demonstrate expressiveness, we \paraphrase{present} a program called a
\emph{module loader} (\cref{module-loader}), written in ILD and making use of
callsite-triggered macros, that is in turn able to execute programs that are
decomposed into \emph{modules} such as the following example:

\begin{ildcode}
(module
  (doc "this module calculates 5!")
  (exports main)
  (imports
    (host (macroexpand <= * +))
    ("core/prelude.ild" (if))
    ("core/module-utils.ild" (fn)))
  (defs
    (!fn main () (fact 5))

    (!fn fact (x)
      (!if (<= x 0)
        1
        (* x (fact (+ x -1)))))))
\end{ildcode}
\label{example-module-fact}
\lstcaption{Example module}

In further research, we aim to also meet the
\clink{c-performance}{performance} constraint by employing partial evaluation
as an optimisation step, which is a technique
known\cite{hudak}\cite{anydsl} to give good results for reducing the
overhead incurred by metaprogramming.
